% =====================================================================
% 第九讲 (Exposé IX) 平展态射的下降。基本群的应用
% (Descente des morphismes étales. Application au groupe fondamental)
% A. Grothendieck
% =====================================================================
\chapter{平展态射的下降。基本群的应用}

\section{平展态射的回顾}

我们在这里回顾平展态射的 (在 I 中发展的) 将要用到的性质，借机从理论中消除多余的诺特假设. 读者会注意，即使只关心诺特概形，下降技巧也导致引入非诺特概形 (例如 $\mathrm{Spec}(\widehat{A} \otimes_A \widehat{A})$，其中 $A$ 是诺特局部环)，而为能应用纤维范畴的语言，必须定义平展态射等概念时不引入诺特限制. 不愿验证或承认下列断言在无诺特假设下成立的读者，可以满足于在 I 的诺特假设下承认它们，只要在随后各节的断言中引入同样的诺特假设，并对推理中出现的非诺特概形使用下面的\cref{def:IX.1.1}.

\begin{definition}\label{def:IX.1.1}
设 $f \colon X \to S$ 是预概形态射，$x$ 是 $X$ 的点. 称 $f$ 在 $x$ 处平展，或 $X$ 在 $x$ 处平展于 $S$ 上，若存在 $s = f(x)$ 的仿射开邻域 $U$，$x$ 的位于 $U$ 上方的仿射开邻域 $V$，仿射诺特概形 $U_0$，平展 (I) 且仿射的 $U_0$-概形 $V_0$，态射 $U \to U_0$，和 $U$-同构
\[
V \xrightarrow{\sim} V_0 \times_{U_0} U.
\]
\end{definition}

注意当 $S$ 局部诺特时，这个术语与 loc. cit. 一致. 类似地称 $f$ 平展，或 $X$ 平展于 $S$ 上，若 $f$ 在 $X$ 的每个点 $x$ 处平展. 在这些定义下，下列命题毫无困难地归结到诺特情形，在 I 第 4，5，7 号中证明了. 细节读者可参见 EGA IV\footnote{精确地说，EGA IV 17，18.}.

\begin{remark}\label{rem:IX.1.2}
若 $f$ 在 $x$ 处平展，则 $f$ "在 $x$ 处有限表现" (\cref{cor:VIII.3.5})，$x$ 在纤维 $f^{-1}(s)$ 中的局部环是 $k(s)$ 的有限可分扩张，最后 $f$ 在 $x$ 处平坦. 可以证明反过来也成立，故\cref{def:IX.1.1} 与 $S$ 局部诺特的情形相同，只需把"在 $x$ 处有限型"换成"在 $x$ 处有限表现". 由于这个结果证明微妙，我们这里不想给出平展态射概念的这一定义，它不便直接证明下面要跟随的性质.
\end{remark}

先注意平凡地有：

\begin{proposition}\label{prop:IX.1.3}
若 $f \colon X \to S$ 平展，则由基变换 $S' \to S$ 得到的任何态射 $f' \colon X' \to S'$ 也平展.
\end{proposition}

故可以说平展态射构成 $\mathbf{Sch}$ 中箭的范畴的纤维子范畴 (参见 VI 11 a)). 本讲的对象是研究这个 $\mathbf{Sch}$ 上纤维范畴的正合性质.

\begin{proposition}\label{prop:IX.1.4}
设 $f \colon X \to S$ 是预概形态射. 为使它是开浸入，必要且充分条件是它平展且根式.
\end{proposition}

参见 I 5.1. 由此得出，若 $X$ 在 $S$ 上平展，$X$ 在 $S$ 上的每个截面是开浸入，再用 \ref{prop:IX.1.4}，得到：

\begin{corollary}\label{cor:IX.1.5}
设 $X$ 是平展 $S$-预概形. 则 $X$ 在 $S$ 上的截面集与 $X$ 的使结构态射诱导的态射 $\Gamma \to S$ 根式且满射的开部分 $\Gamma$ 集一一对应.
\end{corollary}

而且若 $X$ 在 $S$ 上分离，$\Gamma$ 将是 $X$ 的同时开且闭的部分，但无关紧要. — 作显然的基变换，可以把 \ref{cor:IX.1.5} 写成表面上更一般的形式：

\begin{corollary}\label{cor:IX.1.6}
设 $X$ 和 $Y$ 是两个 $S$-预概形，$Y$ 在 $S$ 上平展. 则映射 $f \mapsto \Gamma_f$ (把 $X$ 到 $Y$ 的每个 $S$-态射 $f$ 配给 $X \times_S Y$ 的 $f$ 的图的部分) 是从 $\mathrm{Hom}_S(X, Y)$ 到 $X \times_S Y$ 的使 $\mathrm{pr}_1$ 诱导的态射 $\Gamma \to X$ 根式且满射的开部分 $\Gamma$ 集的双射.
\end{corollary}

\begin{proposition}\label{prop:IX.1.7}
设 $S_0$ 是由拟凝聚幂零理想 (i.e. 使 $S_0$ 与 $S$ 有同样底集的理想) 定义的 $S$ 的子预概形. 则函子 $X \mapsto X \times_S S_0$ (从 $S$ 上平展预概形范畴到 $S_0$ 上平展预概形范畴) 是范畴等价.
\end{proposition}

该函子忠实全的事实是 \ref{cor:IX.1.6} 的直接推论. 本质满射的事实包含在 I 8.3 中. 注意在前面的等价中，$X$ 有限型 i.e. 在 $S$ 上拟有限，(resp. 有限 i.e. 是 $S$ 的平展覆盖)，当且仅当 $X_0$ 在 $S_0$ 上满足类似条件；对分离条件同样评注. 这些事实直接，也包含在下面的 \ref{prop:IX.2.4} 中.

\begin{corollary}\label{cor:IX.1.8}
设 $A$ 是剩余域 $k$ 的完备诺特局部环. 则函子 $B \mapsto B \otimes_A k$ 是 $A$ 上有限平展代数范畴与 $k$ 上有限平展代数范畴 (i.e. 由有限个 $k$ 的有限可分扩张组成的代数) 的等价.
\end{corollary}

\begin{proposition}\label{prop:IX.1.9}
为使 $X$ 是 $S$ 的平展覆盖 i.e. 在 $S$ 上有限且平展，必要且充分条件是 $X$ 与 $S$ 上某个代数 $\mathcal{A}$ (它是有限型局部自由模) 的谱 $S$-同构，且对每个 $s \in S$，$\mathcal{A}_s \otimes_{\mathcal{O}_s} k(s)$ 是 $k(s)$ 上的可分代数，从而在此情形是 $k(s)$ 的有限可分扩张的直积.
\end{proposition}

最后下列结果较不初等，它是 I 8.4 与代数几何中层存在定理 (EGA III 5；参见 \cite{exp09-1} 定理 3) 的合取.

\begin{theorem}\label{thm:IX.1.10}
设 $S$ 是完备诺特局部环的谱，$X$ 是 $S$-概形紧合，$X_0$ 是 $X$ 在 $S$ 的闭点处的纤维 (故 $X_0$ 是 $X$ 的闭子概形). 则限制函子 $X' \mapsto X' \times_X X_0$ 是 $X$ 的平展覆盖范畴与 $X_0$ 的平展覆盖范畴的等价.
\end{theorem}

\section{浸没态射与万有浸没态射}

\begin{definition}\label{def:IX.2.1}
预概形态射 $g \colon S' \to S$ 称为浸没\index[term]{jinmeitaishe@浸没态射}，若它满射，且使 $S$ 成为 $S'$ 的商拓扑空间 (i.e. $S$ 的部分 $U$ 使 $f^{-1}(U)$ 开，则 $U$ 开). 称 $f$ 万有浸没\index[term]{wanyoujinmeitaishe@万有浸没态射}，若对任何态射 $T \to S$，由 $f$ 经基变换得到的态射 $f' \colon T' = S' \times_S T \to T$ 浸没.
\end{definition}

两个浸没 (resp. 万有浸没) 态射的复合浸没 (resp. 万有浸没)，万有浸没态射的基变换给出万有浸没态射 (已为此做了所需的调整)，这是直接的. 若 $fg$ 浸没 (resp. 万有浸没)，则 $f$ 浸没.

\begin{example}\label{exa:IX.2.2}
a) 满射且开的或闭的态射浸没，故万有闭满射或万有开满射态射万有浸没. 例如紧合满射态射万有浸没. 另一方面，忠实平坦拟紧态射万有浸没 (\cref{cor:VIII.4.3}). 这将是对我们最重要的两个情形.
\end{example}

可以对浸没或万有浸没态射 $g \colon S' \to S$ 应用 VIII 4.3 的推理，特别得到：

\begin{proposition}\label{prop:IX.2.3}
假设 $g \colon S' \to S$ 浸没. 则下列映射图式正合：
\[
\mathrm{Ouv}(S) \longrightarrow \mathrm{Ouv}(S') \rightrightarrows \mathrm{Ouv}(S''),
\]
其中 $S'' = S' \times_S S'$，$\mathrm{Ouv}(X)$ 表示预概形 $X$ 的开部分集.
\end{proposition}

\begin{proposition}\label{prop:IX.2.4}
设 $g \colon S' \to S$ 是万有浸没态射，$f \colon X \to Y$ 是 $S$-态射，$f' \colon X' \to Y'$ 是经基变换得到的 $S'$-态射. 为使 $f$ 开 (resp. 闭)，只需 $f'$ 如此. 为使 $f$ 万有开 (resp. 万有闭，resp. 分离)，必要且充分条件是 $f'$ 如此. 若还假设 $g$ 拟紧，$f$ 局部有限型，为使 $f$ 紧合，必要且充分条件是 $f'$ 如此.
\end{proposition}

对最后一点，注意若 $f'$ 紧合故拟紧，则由 VIII 3.3 $f$ 拟紧，故有限型，因为它局部有限型. 另一方面由前面所述它分离且万有闭，故紧合.

\begin{proposition}\label{prop:IX.2.5}
设 $S'$ 是完备诺特局部环的谱 $S$ 上有限型的预概形，假设 $S$ 的闭点 $s$ 的纤维有限，故该纤维诸点 $s'$ 在 $S'$ 中的局部环在 $A = \mathcal{O}_s$ 上有限. 设 $S''$ 是所论诸 $\mathcal{O}_{S',s'}$ 的谱的和概形，看作有限 $S$-概形. 为使 $g \colon S' \to S$ 万有浸没，必要且充分条件是结构态射 $S'' \to S'$ 满射.
\end{proposition}

由于有自然 $S$-态射 $S'' \to S'$，且有限满射态射由 \ref{exa:IX.2.2} 万有浸没，所述条件充分. 证明若 $S'' \to S'$ 不满射，则 $g$ 不万有浸没. 事实上，设 $t$ 是不在 $S''$ 的像中的 $S$ 的点；则存在 $S$-概形 $T$，为离散赋值环的谱，它在 $S$ 中的像是 $\{s, t\}$. 注意 $S''$ 在 $S'$ 中的像开，因为态射 $S'' \to S'$ 局部同构，另一方面该像包含 $S'_s$，且与 $S'_t$ 不相交. 由此得出这最后一个在 $T' = S' \times_S T$ 中的原像开，且等于 $T$ 的闭点的原像. 这表明 $T' \to T$ 不浸没，故 $S' \to S$ 不万有浸没.

\begin{remark}\label{rem:IX.2.6}
用判则 IV 6.3 (诺特空间中可构部分开) 容易得到下列赋值判则：有限型态射 $g \colon S' \to S$ ($S$ 局部诺特) 万有浸没，必要且充分条件是：对每个为离散赋值环谱的 $S$-概形 $T$，令 $T' = S' \times_S T$，$T$ 的闭点在 $T'$ 中的原像不开.
\end{remark}

\section{平展预概形态射的下降}

\begin{proposition}\label{prop:IX.3.1}
设 $g \colon S' \to S$ 是预概形满射态射，$X$ 和 $Y$ 是 $S$ 上的两个预概形，$X'$，$Y'$ 是它们在 $S'$ 上的逆像. 若 $Y$ 在 $S$ 上非分歧，则典范映射
\[
\mathrm{Hom}_S(X, Y) \longrightarrow \mathrm{Hom}_{S'}(X', Y')
\]
单射.
\end{proposition}

事实上，由 \ref{cor:IX.1.6}，$S$-态射 $f \colon X \to Y$ 在已知它的图的底集 $\Gamma$ (它是 $Z = X \times_S Y$ 的部分) 时确定. 由于
\[
Z' = Z \times_S S' = X' \times_{S'} Y' \longrightarrow Z
\]
满射 (因为 $S' \to S$ 满射)，该部分 $\Gamma$ 在已知它在 $X' \times_{S'} Y'$ 中的原像时确定，后者正是 $f'$ 的图的底集. 由此得结论.

$Z$ 的部分 $\Gamma$ 是 $S$-态射 $f \colon X \to Y$ 的图，当且仅当它开，且 $\mathrm{pr}_1$ 诱导的态射 $\Gamma \to Y$ 根式且满射，参见 1. 当第一个性质满足时，第二个性质成立当且仅当 $\Gamma$ 在 $Z'$ 中的原像 $\Gamma'$ 满足同样条件 (\cref{prop:VIII.3.1}). 若还知道 $Z' \to Z$ 浸没 (特别当 $S' \to S$ 万有浸没时如此)，则 $\Gamma$ 开当且仅当 $\Gamma'$ 开. 这样，集合 $\mathrm{Hom}_S(X, Y)$ 于是与 $Z'$ 的下列开部分 $\Gamma'$ 的集一一对应：投影态射 $\mathrm{pr}_1 \colon Z' \to X'$ 根式且满射 (i.e. 对应 $S'$-态射 $f' \colon X' \to Y'$)，且对由 $Z' \to Z$ 定义的等价关系饱和，i.e. 它在 $Z'' = Z' \times_Z Z' = Z \times_S S''$ (其中 $S'' = S' \times_S S'$) 中经两个投影之一的另一个的原像相等. 而后者正是由 $f'$ 经基变换 ($S'' \to S'$ 的两个投影之一另一个) 得到的两个 $S''$-态射 $X'' \to Y''$ 的图. 于是得到：

\begin{proposition}\label{prop:IX.3.2}
设 $g \colon S' \to S$ 是预概形万有浸没态射，$S'' = S' \times_S S'$，$X$ 和 $Y$ 是两个 $S$-预概形，$X'$ 和 $Y'$ 是它们在 $S'$ 上的逆像，$X''$，$Y''$ 是它们在 $S''$ 上的逆像. 若 $Y$ 在 $S$ 上平展，则下列典范映射图式正合：
\[
\mathrm{Hom}_S(X, Y) \longrightarrow \mathrm{Hom}_{S'}(X', Y') \rightrightarrows \mathrm{Hom}_{S''}(X'', Y'').
\]
\end{proposition}

取 $X$ 和 $Y$ 在 $S$ 上平展，得到下列断言，它反过来给出 \ref{prop:IX.3.2} (甚至限于 $X = S$，此时总能由基变换 $X \to S$ 归结到 \ref{prop:IX.3.2})：

\begin{corollary}\label{cor:IX.3.3}
预概形的万有浸没态射对互相平展的预概形的纤维范畴是下降态射.
\end{corollary}

我也不知道它是否必然对所述纤维范畴是有效下降态射，即使再假设 $S$ 诺特且 $g$ 有限型，并限于平展覆盖. 我们仍将在下一节给出有用的有效性判据.

\begin{corollary}\label{cor:IX.3.4}
设 $g \colon S' \to S$ 是万有浸没态射，其诸纤维 $g^{-1}(s)$ "几何连通\index[term]{jiheliantong@几何连通}"，i.e. 对任何扩张 $K/k(s)$，$g^{-1}(s) \otimes_{k(s)} K$ 连通. 则 $S$ 连通时 $S'$ 连通. 由 $g$ 定义的从 $S$ 上平展预概形范畴到 $S'$ 上平展预概形范畴的函子忠实全.
\end{corollary}

$S'$ 的同时开且闭的部分对由 $g$ 定义的集合等价关系饱和，因为诸纤维连通，故它是 $S$ 的某个部分的原像，后者由于 $g$ 浸没必然开且闭. 故 $S$ 连通时 $S'$ 连通. 这也蕴含下列结果：两个万有连通纤维态射的复合 $fg$ ($f$ 万有浸没) 有万有连通纤维；若 $S$ 上的 $S'_1$ 和 $S'_2$ 有万有连通纤维，则 $S'_1 \times_S S'_2$ 亦然. 特别在 \ref{cor:IX.3.4} 的条件下，$S''$ 在 $S$ 上有万有连通纤维. 设 $X$ 和 $Y$ 在 $S$ 上平展，设 $u'$ 是 $X'$ 到 $Y'$ 的 $S'$-态射，证明它与下降数据相容 (这由 \ref{cor:IX.3.3} 给出所求结论). 设 $u''_1$ 和 $u''_2$ 是由 $u'$ 得到的两个 $S''$-态射 $X'' \to Y''$. $S''$ 的 $u''_1$ 和 $u''_2$ 的一致点子预概形是诱导开子预概形，纤维性地闭，它是 $Y''$ 在 $S''$ 上的对角预概形的原像\footnote{注意 $S'$ 在 $S$ 上的诸纤维分离!}. 故它是 $S$ 的某部分的原像. 由于它包含 $S''$ 中的对角线，它等于 $S''$，故 $u''_1 = u''_2$，证毕.

\section{平展预概形的下降：有效性判据}

\begin{proposition}\label{prop:IX.4.1}
设 $g \colon S' \to S$ 是忠实平坦拟紧态射. 则 $g$ 对互相平展、分离且有限型的预概形的纤维范畴是有效下降态射.
\end{proposition}

这确实是所述纤维范畴的下降态射，由 \ref{cor:IX.3.3} 或 VIII 5.2 任选其一. 剩下证明：若 $X'$ 在 $S'$ 上平展、分离且有限型，配备相对 $g \colon S' \to S$ 的下降数据，则该数据在所述纤维范畴中有效. 容易看到，若 $X$ 是 $S$ 上的预概形，则它在 $S$ 上平展当且仅当它在 $S'$ 上平展 (由\cref{def:IX.1.1} 和 loc. cit. 3.6). 故它在 $S$ 上平展、分离且有限型，当且仅当 $X'$ 在 $S'$ 上如此，参见例如 2.4. 故只需保证 $X$ 上下降数据对 $\mathbf{Sch}$ 的箭的纤维范畴有效. 而由 VIII 6.2 和 \ref{cor:IX.6.6}，$X'$ 在 $S'$ 上拟仿射. 于是可用 VIII 7.9 得出结论. 读者还会注意，若限于互相平展且有限的预概形，证明要求更少，因为可直接援引 VIII 2.1.

\begin{corollary}\label{cor:IX.4.2}
设 $g \colon S' \to S$ 是万有浸没态射，$X'$ 是 $S'$ 上配备相对 $g$ 的下降数据的平展分离有限型 $S'$-预概形，$S_1 \to S$ 是忠实平坦拟紧态射，$S'_1$ 和 $X'_1$ 由 $S'$ 和 $X'$ 经基变换得到，故 $S'_1 \to S_1$ 万有浸没，$X'_1$ 在 $S'_1$ 上平展分离且有限型，配备相对 $g_1 \colon S'_1 \to S_1$ 的下降数据. 为使 $X'$ 上的下降数据有效，必要且充分条件是 $X'_1$ 上的下降数据有效.
\end{corollary}

这由范畴中的下降理论 \cite{exp09-D} 得出，考虑到 \ref{prop:IX.4.1} 和 3.3.

类似地证明：

\begin{corollary}\label{cor:IX.4.3}
设 $g \colon S' \to S$ 是万有浸没态射，$X'$ 是 $S'$ 上配备相对 $g$ 的下降数据的平展 $S'$-预概形，$(S_i)$ 是 $S$ 的开覆盖. 为使下降数据有效，必要且充分条件是：对每个 $i$，相对态射 $g_i \colon S'_i = S' \times_S S_i \to S_i$ 在 $X'_i = X' \times_S S_i$ 上对应的下降数据有效.
\end{corollary}

最后一个结果导致给出局部有效性判据：

\begin{proposition}\label{prop:IX.4.4}
设 $g \colon S' \to S$ 是有限表现 (\cref{rem:VIII.3.6}) 且万有浸没的态射，$X'$ 是 $S'$ 上平展有限表现的预概形，配备相对 $g$ 的下降数据，最后 $a$ 是 $S$ 的点. 为使存在 $a$ 的开邻域 $U$，使相对态射
\[
g_U \colon S'_U = S' \times_S U \longrightarrow S_U = U
\]
在 $X'_U = X' \times_S U$ 上对应的下降数据有效，必要且充分条件是相对态射
\[
g_a \colon S'_a = S' \times_S \mathrm{Spec}(\mathcal{O}_a) \longrightarrow S_a = \mathrm{Spec}(\mathcal{O}_a)
\]
在 $X'_a = X' \times_S \mathrm{Spec}(\mathcal{O}_a)$ 上对应的下降数据有效.
\end{proposition}

必要性平凡，证明充分性. 于是有 $S_a$ 上有限型平展预概形 $X_a$ 和与下降数据相容的同构
\[
(\ast) \quad X'_a \xrightarrow{\sim} X_a \times_{S_a} S'_a.
\]
按关于环的归纳极限 (这里是环 $A_f$，其中 $A$ 是 $a$ 的仿射开邻域的环，$f$ 遍历不在对应 $a$ 的素理想中的元素) 上定义的预概形的容易的一般推论，可以找到 $a$ 的开邻域 $U$，$U = S_U$ 上有限型平展预概形 $X_U$，和 $S_a$-同构 $X_a \xrightarrow{\sim} X_U \times_{S_U} S_a$. 而且，取 $U$ 足够小，可假设同构 $(\ast)$ 来自同构：
\[
X'_U \xrightarrow{\sim} X_U \times_{S_U} S'_U;
\]
后者可能不与下降数据相容，但只要缩小 $U$，它将与下降数据相容. 这完成证明.

\begin{corollary}\label{cor:IX.4.5}
在 \ref{prop:IX.4.4} 的条件下，为使 $X'$ 上的下降数据有效，必要且充分条件是：对每个 $a \in S$，相对态射 $S'_a = S' \times_S \mathrm{Spec}(\mathcal{O}_a) \to \mathrm{Spec}(\mathcal{O}_a)$ 在 $X'_a$ 上对应的下降数据有效. 当 $S$ 局部诺特且 $X'$ 在 $S'$ 上分离时，在前面判则中还可把 $\mathcal{O}_a$ 换成它的完备化.
\end{corollary}

第一个断言由 \ref{prop:IX.4.4} 和 \ref{cor:IX.4.3} 得出，第二个断言则是 \ref{cor:IX.4.2} 的推论. 再用 \ref{cor:IX.4.2} 及事实：对每个诺特局部环 $A$，可以找到完备诺特局部环 $B$ 和局部同态 $A \to B$，使 $B$ 在 $A$ 上平坦且 $B/\mathfrak{m}B$ 是 $A$ 的剩余域 $k = A/\mathfrak{m}$ 的给定扩张，得到：

\begin{corollary}\label{cor:IX.4.6}
在 \ref{prop:IX.4.4} 的条件下，再假设 $X'$ 在 $S'$ 上分离，$S$ 局部诺特. 为使 $X'$ 上的下降数据有效，必要且充分条件是：对每个为剩余域代数闭的完备局部环谱的 $S$-概形 $S_1$，相对态射 $g_1 \colon S'_1 \to S_1$ 在 $X'_1 = X' \times_S S_1$ 上对应的下降数据有效.
\end{corollary}

\begin{theorem}\label{thm:IX.4.7}
设 $g \colon S' \to S$ 是有限满射有限表现态射 (若 $S$ 局部诺特，这最后一个假设是其他假设的推论)\footnote{可以证明实际上只需 $g$ 是整态射，用 EGA IV 8 风格的极限过程归结到正文的情形.}. 则 $g$ 对互相平展、分离、有限型的预概形的纤维范畴是有效下降态射.
\end{theorem}

需证明：若 $X'$ 在 $S'$ 上平展、分离、有限型，配备相对 $g$ 的下降数据，则该数据有效. 用 \ref{cor:IX.4.3} 容易归结到 $S$ 诺特的情形. 由 \ref{cor:IX.4.5}，可假设 $S$ 是诺特局部环的谱，从而尤其
\[
\dim S = n < +\infty.
\]
对 $\dim S = n$ 作递归，$n < 0$ 时断言平凡. 假设 $n \geqslant 0$ 且定理对维数 $n' < n$ 已证. 由 \ref{cor:IX.4.6} 归结到 $S$ 是完备局部环的谱的情形，故 $S'$ 是完备局部环的谱的有限并. 故有
\[
X' = X'_1 \sqcup X'_2
\]
其中 $X'_1$ 在 $S'$ 上有限，$X'_2$ 在 $S'$ 的任一个闭点上方都没有点. 考虑对应下降数据的态射
\[
q_1, q_2 \colon X'' \rightrightarrows X',
\]
它与 $p_1, p_2 \colon S'' \rightrightarrows S'$ 相容. 立即看到
\[
X'' = q_i^{-1}(X'_1) \sqcup q_i^{-1}(X'_2) \quad i = 1, 2
\]
是 $X''$ 在 $S''$ 上的对应典范分解，这蕴含 $q_1^{-1}(X'_1) = q_2^{-1}(X'_1)$，从而 $X'_1$ 和 $X'_2$ 配备诱导下降数据. 设 $T$ 是 $S$ 的闭点的余集的开集，故 $T' = S' \times_S T$ 是 $S'$ 的闭点集余集的部分，而整个位于 $T'$ 上方的 $X'_2$ 配备相对 $g$ 诱导的态射 $T' \to T$ 的下降数据. 由于后者有限满射，且 $\dim T < \dim S = n$，由递归假设该下降数据有效. 故只需证明 $X'_1$ 上的下降数据有效，现在可假设 $X'$ 在 $S'$ 上平展且有限. (N.B. 若在断言 \ref{thm:IX.4.7} 中限于平展覆盖，递归推理不必要). 设 $S_0$ 是 $A$ 的剩余域的谱，设 $S'_0 = S' \times_S S_0$，同样由 $S'$ 在 $S$ 上的纤维平方和立方 $S''$ 和 $S'''$ 定义 $S''_0$，$S'''_0$. 由 \ref{cor:IX.1.8}，态射 $S_0 \to S$，$S'_0 \to S'$ 等对平展覆盖范畴 (分别 $S$ 和 $S_0$ 的，$S'$ 和 $S'_0$ 的，等等) 诱导等价. 由范畴中下降理论的推论 \cite{exp09-D}，得出 $g \colon S' \to S$ 对平展覆盖纤维范畴是有效下降态射，当且仅当 $g_0 \colon S'_0 \to S_0$ 如此. 而确实如此，例如作为 \ref{prop:IX.4.1} 的特殊情形. 这完成 \ref{thm:IX.4.7} 的证明.

\begin{corollary}\label{cor:IX.4.8}
假设 $S$ 局部诺特，\ref{thm:IX.4.7} 的结论在只假设 $S' \to S$ 万有浸没、有限型且拟有限时仍成立.
\end{corollary}

由 \ref{cor:IX.4.6}，可假设 $S$ 是完备诺特局部环的谱. 则由 \ref{prop:IX.2.5}，存在有限满射态射 $S_1 \to S$ 和 $S$-态射 $S_1 \to S'$. 由于 $S_1 \to S$ 由 \ref{thm:IX.4.7} 对所述纤维范畴是万有严格下降态射，且 $S' \to S$ 对它万有下降，\ref{cor:IX.4.8} 由一般推论 \cite{exp09-D} 得出.

\begin{corollary}\label{cor:IX.4.9}
设 $g \colon S' \to S$ 是有限型、满射且万有开的态射，$S$ 局部诺特. 则 $g$ 对互相平展、分离且有限型的预概形的纤维范畴是有效下降态射.
\end{corollary}

如 \ref{thm:IX.4.7} 那样进行，归结到 $S$ 是诺特完备局部环 $A$ 的谱的情形. 设 $A_1$ 是 $A$ 上的有限代数，谱为 $S_1$，使 $S_1 \to S$ 有限满射，故由 \ref{thm:IX.4.7} 对所述纤维范畴是万有有效下降态射. 由一般定理 \cite{exp09-D} 得出，$g$ 对所述纤维范畴是有效下降态射，当且仅当对应态射 $g_1 \colon S'_1 = S' \times_S S_1 \to S_1$ 如此. 由于后者满足与 $g$ 同样的假设，归结到对 $S_1$ 代替 $S$ 证明 4.9. 先取 $A_1$ 为诸 $A/\mathfrak{p}_i$ 的直积 (其中 $\mathfrak{p}_i$ 遍历 $A$ 的极小素理想)，归结到 $A$ 整的情形. 此时证明\footnote{参见 EGA IV 14.3.13 和 14.5.4.}存在 $S'$ 的整子概形 $S_1$，在 $S$ 上拟有限且支配 $S$，经过 $S'$ 在 $S$ 的闭点 $y$ 处的纤维的某个点 (因为 $S'$ 在 $S$ (整局部诺特) 上万有开有限型，且 $S'_y \neq \varnothing$). 由于 $A$ 完备，$S_1$ 在 $S$ 上有限，且因它支配 $S$，态射 $S_1 \to S$ 满射. 再次把 $S$ 换成 $S_1$，归结到 $S'$ 在 $S$ 上有截面的情形，此时断言平凡.

\begin{theorem}\label{thm:IX.4.10}
设 $g \colon S' \to S$ 是有限根式满射有限表现态射 (若 $S$ 局部诺特，这最后一个条件多余\footnote{甚至只需 $g$ 整、根式且满射，如容易归结到正文情形所见，风格如 EGA IV 8，参见 SGA 4 VIII 1.1.}). 则逆像函子诱导 $S$ 上平展预概形范畴与 $S'$ 上平展预概形范畴的等价.
\end{theorem}

由于 $S'$ 在 $S' \times_S S'$ 和 $S' \times_S S' \times_S S'$ 中的对角态射是满射浸入，故由 \ref{prop:IX.1.9} 诱导 (分别 $S' \times_S S'$ 和 $S' \times_S S' \times_S S'$ 上平展预概形范畴与 $S'$ 上平展预概形范畴的) 等价，由下降推论 \cite{exp09-D} 得出每个在 $S'$ 上平展的 $X'$ 配备相对 $g \colon S' \to S$ 的下降数据，且只有一个. 故 \ref{cor:IX.3.3} 蕴含经 $g$ 的逆像函子 (从 $S$ 上平展预概形范畴到 $S'$ 上平展预概形范畴) 忠实全. 剩下证明它本质满射，i.e. 每个在 $S'$ 上平展的 $X'$ 同构于某个在 $S$ 上平展的 $X$ 的逆像. 问题显然在 $S$ 和 $X'$ 上局部，可假设 $S$，$S'$，$X'$ 仿射. 但此时 $X'$ 在 $S'$ 上分离且有限型，可应用有效性判据 4.7.

\begin{corollary}\label{cor:IX.4.11}
把对 $g$ 的假设换成：$g$ 忠实平坦、拟紧且根式，则 \ref{cor:IX.4.9} 的结论仍然成立.
\end{corollary}

同样的证明，用 \ref{prop:IX.4.1} 代替 4.7.

注意 \ref{thm:IX.4.7} 的证明"初等"，因为它不用紧合态射的有限性与比较定理 (EGA III 3，4，5). 对下列结果不再如此：

\begin{theorem}\label{thm:IX.4.12}
设 $g \colon S' \to S$ 是紧合、满射、有限表现 (若 $S$ 局部诺特，这最后一个假设是第一个的推论) 的态射. 则 $g$ 对预概形的平展覆盖纤维范畴是有效下降态射.
\end{theorem}

由 \ref{cor:IX.3.3} 和 \ref{exa:IX.2.2}，归结到证明：对 $S'$ 上的每个平展覆盖 $X'$，配备相对 $g \colon S' \to S$ 的下降数据，该下降数据有效. 用 \ref{cor:IX.4.3} 容易归结到 $S$ 诺特的情形，用 \ref{cor:IX.4.6} 可假设 $S$ 是完备诺特局部环 $A$ 的谱. 如常引入 $S''$ 和 $S'''$，设 $S_0$ 是 $A$ 的剩余域的谱，设 $S'_0$，$S''_0$，$S'''_0$ 由 $S'$，$S''$，$S'''$ 经基变换 $S_0 \to S$ 得到，i.e. 是 $S'$，$S''$，$S'''$ 在 $S$ 的闭点处的纤维. 由 \ref{thm:IX.1.10}，态射 $S_0 \to S$，$S'_0 \to S'$ 等对目标概形与源概形的平展覆盖范畴诱导等价. 从而 $g \colon S' \to S$ 对预概形的平展覆盖纤维范畴是严格下降态射，当且仅当 $g_0 \colon S'_0 \to S_0$ 如此，而由 \ref{prop:IX.4.1} 确实如此. 这完成 \ref{thm:IX.4.12} 的证明. (在这个推理中，只需要 \ref{thm:IX.1.10} 中所述函子忠实全的事实，它不用代数几何中凝聚层存在定理.)

\section{基本群术语的翻译}

设
\[
g \colon S' \longrightarrow S
\]
是对预概形的平展覆盖纤维范畴的有效下降态射，例如紧合、满射、有限表现态射 (\cref{thm:IX.4.12})，或忠实平坦拟紧态射. 如常引入 $S''$，$S'''$，并用 $\mathcal{C}$，$\mathcal{C}'$，$\mathcal{C}''$，$\mathcal{C}'''$ 表示 $S$，$S'$，$S''$，$S'''$ 的平展覆盖范畴，有范畴的 2-正合图式
\[
(\ast) \quad \mathcal{C} \xrightarrow{\;p^{*}\;} \mathcal{C}' \xrightarrow{\;p_1^{*}, p_2^{*}\;} \mathcal{C}'' \xrightarrow{\;p_{21}^{*}, p_{32}^{*}, p_{31}^{*}\;} \mathcal{C}'''
\]
(其中第二个和第三个箭分别表示两个箭 $p_1^{*}, p_2^{*}$ 和三个箭 $p_{21}^{*}, p_{32}^{*}, p_{31}^{*}$) 对应图式
\[
S \longleftarrow S' \xleftarrow{p_1, p_2} S'' \xleftarrow{p_{21}, p_{32}, p_{31}} S'''.
\]
假设预概形 $S$，$S'$，$S''$，$S'''$ 是连通预概形的无交和，特别当它们是局部连通预概形时如此，更不用说局部诺特 (例如 $S'$ 在局部诺特 $S$ 上有限型). 则 $(\ast)$ 中的范畴 $\mathcal{C}$，$\mathcal{C}'$ ... 是 V 9 的多 Galois 范畴，各自由一族拓扑紧全不连通的群描述，即预概形 $S$，$S'$，$S''$，$S'''$ 的连通分量的基本群. 为简单假设 $S$ 连通，将给出计算它的基本群的方法，用由 $\mathcal{C}'$，$\mathcal{C}''$，$\mathcal{C}'''$ 组成、借助表达这些范畴的基本群适当明确的纤维范畴. 读者会注意，勾勒的方法实际上在多 Galois 范畴的一般框架 (它们不必来自给定的预概形 $S$，$S'$，$S''$，$S'''$) 中有效. 它也是计算拓扑空间 $S$ (闭子空间 $S_i$ 的局部有限并，或开子空间 $S_i$ 的任意并) 基本群的熟知方法的类似物，用 $S_i$ 分量和 $S_i \cap S_j$ 分量的基本群. 当然，预概形框架中的类似情形确实落入下降的一般框架，只需引入 $S_i$ 的和预概形 $S'$ 和典范态射 $g \colon S' \to S$.

令
\[
E' = \pi_0(S'), \quad E'' = \pi_0(S''), \quad E''' = \pi_0(S'''),
\]
其中 $\pi_0$ 表示"连通分量集"函子. 由于 $S'$ 在 $S$ 上的纤维积构成 $\mathbf{Sch}$ 的单纯对象，它被函子 $\pi_0$ 变换成 $E'$，$E''$，$E'''$ 是维数 0，1，2 的分量的单纯集. 将用到单纯映射
\[
q_i = \pi_0(p_i), \quad (i = 1, 2), \quad \text{和} \quad q_{ij} = \pi_0(p_{ij}), \quad (i, j) = (2, 1), (3, 2), (3, 1),
\]
在图式中显示：
\[
(1) \quad E' \xleftarrow{\;q_1, q_2\;} E'' \xleftarrow{\;q_{21}, q_{32}, q_{31}\;} E'''.
\]
$E'$ 的对象用带一个撇的记号，如 $s'$，$E''$ resp. $E'''$ 的对象用带两个 resp. 三个撇的记号. $S$ 连通的事实翻译为 $\pi_0(K) = 0$，其中 $K$ 是由 $g \colon S' \to S$ 定义的单纯集，或也翻译为：映射对 $(q_1, q_2)$ 在 $E'$ 中生成的等价关系传递.

一次选定 $E'$ 的元素 $s'_0$，并对每个 $s' \in E'$ 选 $E''$ 的元素 $s'$ 使\footnote{注意隐式承认存在的元素 $s'$ 并非在所有情形存在. 故如所述陈述的\cref{thm:IX.5.1} 不适用于所有情形. 但参照所写的文本不难修改该定理的陈述，使它给出在所有情形适用的计算方法. 特别，该定理的诸推论原样成立.}
\[
q_1(s') = s'_0, \quad q_2(s') = s',
\]
从而显示 $S$ 的连通性. 对每个 $s' \in E'$，在 $S'$ 的连通分量 $s'$ 中选几何点 $s'$；该点实际上通过 $C'$ 上对应的纤维函子 $F'_{s'}$ 出现. 该函子的自同构群，i.e. $S'$ 在 $s'$ 处的基本群，记为 $\pi_{s'}$. 类似地选 $s''$ 和 $s'''$，从而选函子 $F''_{s''}$ 和 $F'''_{s'''}$，得基本群 $\pi_{s''}$ 和 $\pi_{s'''}$. 这样
\[
\pi_{s'} = \pi_1(S', s'), \quad \pi_{s''} = \pi_1(S'', s''), \quad \pi_{s'''} = \pi_1(S''', s''').
\]
对每个 $s'' \in E''$，$p_1(s'')$ 与 $q_1(s'')$ 在同一个连通分量中，故存在函子同构 $F''_{s''} \circ p_1^{*} \xrightarrow{\sim} F'_{q_1(s'')}$ (i.e. 从 $p_1(s'')$ 到 $q_1(s'')$ 的"道路类"). 对 $q_2$ 和诸 $q_{ij}$ 重复此评注. 选定所有这些道路类：
\[
F''_{s''} \circ p_i^{*} \xrightarrow{\sim} F'_{q_i(s'')}, \quad F'''_{s'''} \circ p_{ij}^{*} \xrightarrow{\sim} F''_{q_{ij}(s''')},
\]
(对 $i = 1, 2$ 和 $(i, j) = (2, 1), (3, 2), (3, 1)$). 特别由此得到群同态：
\[
(2) \quad q_i^{s''} \colon \pi_{s''} \longrightarrow \pi_{q_i(s'')}, \quad q_{ij}^{s'''} \colon \pi_{s'''} \longrightarrow \pi_{q_{ij}(s''')},
\]
($i$ 和 $(i, j)$ 取同样值). 最后回忆在纤维 $\mathcal{C}'$，$\mathcal{C}''$，$\mathcal{C}'''$ 的劈开范畴结构中还有函子同构：
\[
p_{21}^{*} p_1^{*} \xrightarrow{\sim} p_{31}^{*} p_1^{*}, \quad p_{21}^{*} p_2^{*} \xrightarrow{\sim} p_{32}^{*} p_1^{*}, \quad p_{31}^{*} p_2^{*} \xrightarrow{\sim} p_{32}^{*} p_2^{*},
\]
它们由两个成员分别与 $u_i^{*}$ ($i = 1, 2, 3$，其中 $u_i$ 是 $S'''$ 到 $S'$ 的三个投影) 的同构得到. 明确这些数据时，对每个 $s'''$ 找到完全确定的元素
\[
(3) \quad a_i^{s'''} \in \pi_{v_i(s''')},
\]
(其中 $v_i$，$i = 1, 2, 3$ 是由 $v_i = \pi_0(u_i)$ 定义的三个映射 $E''' \to E'$)，它们服从条件：
\[
q_1^{s''_1} q_{21}^{s'''} = \mathrm{int}(a_1^{s'''}) q_1^{s''_2} q_{31}^{s'''} \quad (s''_1 = q_{21}(s'''), s''_2 = q_{31}(s''')),
\]
以及两个类似条件，涉及 $a_2$ 和 $a_3$. 读者还会注意，数据 (1)，(2)，(3) 允许在纤维范畴等价的范围内重构所考虑的纤维 $\mathcal{C}'$，$\mathcal{C}''$，$\mathcal{C}'''$ 的纤维范畴. 故原则上应允许在等价范围内重构 $\mathcal{C}$，从而在同构范围内重构它的基本群. 我们实际上将在几何点 $p(s'_0)$ 处确定基本群，i.e. 函子 $F'_{s'_0} \circ p^{*}$ 的自同构群.

注意 $\mathcal{C}'$ 的对象 $X'$ 的给出本质上等价于有限集 $X'_{s'}$ ($s' \in E'$) 的给出，其中 $\pi_{s'}$ 连续作用. 这样的对象上的粘合数据于是等价于对每个 $s'' \in E''$ 给一个双射：
\[
\varphi_{s''} \colon X'_{q_1(s'')} \xrightarrow{\sim} X'_{q_2(s'')}
\]
与 $\pi_{s''}$ 的作用相容，后者借助同态 $q_i^{s''} \colon \pi_{s''} \to \pi_{q_i(s'')}$ 作用于两个成员之一另一个. 先取 $s''$ 为 $s'$ 形式的元素，看到这样的数据定义双射
\[
\psi_{s'} \colon X'_{s'_0} = F'_0(X') \xrightarrow{\sim} X'_{s'}
\]
这使得可以把诸 $X'_{s'}$ 等同于同一个集合 $F'_0(X') = X'_{s'_0}$，所有群 $\pi_{s'}$ 此后在其上作用. 这样，双射 $\varphi_{s''}$ 对应双射
\[
g_{s''} \colon F'_0(X') \xrightarrow{\sim} F'_0(X'),
\]
一方面服从与 $\pi_{s''}$ 的交换关系：
\[
\mathrm{a)} \quad g_{s''} q_1^{s''}(g'') = q_2^{s''}(g'') g_{s''} \quad (s'' \in E'', g'' \in \pi_{s''}),
\]
另一方面服从关系
\[
\mathrm{b)} \quad g_{s'} = g_{s'_0} \quad (s' \in E'),
\]
表达我们把诸 $X_{s'}$ 彼此等同的方式. 明确这样的粘合数据成为下降数据的条件时，找到关系：
\[
\mathrm{c)} \quad a_3^{s'''} g_{q_{31}(s''')} a_1^{s'''} = g_{q_{32}(s''')} a_2^{s'''} g_{q_{21}(s''')} \quad (s''' \in E''').
\]
这给出 $\mathcal{C}'$ 的配备下降数据的对象范畴与下列范畴的等价：有限集 ($\pi_{s'}$ 在其上连续作用)，再配备满足关系 a)，b)，c) 的双射 $g_{s''}$.

设 $G$ 是由群 $\pi_{s'}$ 和新生成元 $g_{s''}$ 服从关系 a)，b)，c) 生成的群，设 $\pi$ 是 $G$ 的使在诸群 $\pi_{s'}$ 中的原像为开子群的有限指数子群的商的投射极限. 也说 $\pi$ 是由诸 $\pi_{s'}$ 和诸 $g_{s''}$ 服从关系 a)，b)，c) 生成的 Galois 型群. 立即看到所考虑的范畴也等价于拓扑群 $\pi$ 在其上连续作用的有限集范畴. 这建立下列断言：

\begin{theorem}\label{thm:IX.5.1}
设 $g \colon S' \to S$ 是对预概形的平展覆盖纤维范畴的有效下降态射 (参见 \ref{cor:IX.4.9} 和 \ref{thm:IX.4.12}) 的预概形态射. 假设 $S$ 连通，且 $S'$，它的纤维平方 $S''$ 和纤维立方 $S'''$ 是连通预概形的和 (例如 $S'$ 在连通局部诺特 $S$ 上有限型时如此). 如上选择：$S'$，$S''$，$S'''$ 的每个连通分量中的几何点，某些道路类，$E'$ 中的元素 $s'_0$，以及对每个 $s' \in E'$ 满足在 $E'$ 中的两个像为 $s'_0$ 和 $s'$ 的 $E''$ 的元素 $s'$. ($E'$，$E''$，$E'''$ 分别表示 $S'$，$S''$，$S'''$ 的连通分量集). 则 $S$ 在 $s'_0$ 的像处的几何点的基本群典范同构于由诸 $\pi_{s'} = \pi_1(S', s')$ ($s' \in E'$) 和生成元 $g_{s''}$ ($s'' \in E''$) 服从上述关系 a)，b)，c) (涉及群 $\pi_{s''} = \pi_1(S'', s'')$ 的元素和前面引入的元素 $a_i^{s'''}$ ($i = 1, 2, 3$，$s''' \in E'''$)) 生成的 Galois 型群.
\end{theorem}

\begin{corollary}\label{cor:IX.5.2}
假设 $S'$ 和 $S''$ 只有有限个连通分量，且 $S'$ 的连通分量的基本群拓扑有限生成. 则 $S$ 的基本群拓扑有限生成.
\end{corollary}

这样，我们以后将证明代数闭域上射影\emph{正规概形}的基本群拓扑有限生成. 用 Chow 引理和代数概形的正规化，由此得出同样结果对代数闭域上任何紧合概形成立.

\begin{corollary}\label{cor:IX.5.3}
假设 $S'$，$S''$，$S'''$ 只有有限个连通分量，$S'$ 的连通分量的基本群拓扑有限表现，$S''$ 的连通分量的基本群拓扑有限生成. 则 $S$ 的基本群拓扑有限表现.
\end{corollary}

注意可以把 \ref{cor:IX.4.9} (限于平展覆盖) 表达为：诺特预概形的有限根式满射态射诱导基本群的同构；形象地说，基本群是预概形的拓扑不变量. 可以用 \ref{thm:IX.5.1} 更一般地明确对预概形操作 (例如预概形沿有限点集的"捏合"，它有简单拓扑意义) 对基本群的作用. 例如得到：

\begin{corollary}\label{cor:IX.5.4}
设 $g \colon S' \to S$ 是有限有限表现态射，$T$ 是 $S$ 的离散部分. 对每个 $s \in S$，设 $n(s)$ 是纤维 $g^{-1}(s)$ 中的"\emph{几何点数}" (它也可明确为 $g^{-1}(s)$ 在 $k(s)$ 上的可分次数，即它的剩余扩张的可分次数的和). 假设对 $s \in S - T$，有 $n(s) = 1$. 对每个 $s \in T$，设 $K_s$ 是 $k(s)$ 的代数闭扩张，$I_s$ 是 $S'$ 的在 $K_s$ 中取值的几何点集 (这是有 $n(s)$ 个元素的集合)，$I'_s$ 是 $I_s$ 去掉选定一个点的余集，最后 $I'$ 是诸 $I'_s$ 的并. 假设 $S'$ 连通. 则 $S$ 的基本群同构于由 $S'$ 的基本群和生成元 $g_i$ ($i \in I'$) 不受任何附加条件生成的 Galois 型群.
\end{corollary}

证明细节留给读者；所得断言只是下述事实在群论语言中的翻译：$S$ 的平展覆盖范畴 $\mathcal{C}$ 与 $S'$ 的平展覆盖 $X'$ 范畴 (对每个 $s \in T$，配备 $X'$ 在 $g^{-1}(s)$ 的诸点在 $K_s$ 中取值的 $n(s)$ 个纤维之间的传递双射系统) 等价. (当然在这个内蕴形式下，不再需要假设 $S'$ 连通).

\begin{example}\label{exa:IX.5.5}
容易证明代数闭域 $k$ 上的有理曲线 $\mathbf{P}^1_k$ 单连通\footnote{参见 XI 1.1.}. 故恰有一个二重点、有 $n$ 个解析分支的完备有理曲线的基本群是由 $n - 1$ 个生成元生成的自由 Galois 型群. 例如在普通二重点的情形，得到基本群 $\widehat{\mathbb{Z}}$，如 (I 11 a)) 中所宣告. 相反，尖点 (它是"几何单支"的点) 的存在对基本群没有影响.
\end{example}

\begin{corollary}\label{cor:IX.5.6}
设 $g \colon S' \to S$ 是纤维几何连通的万有浸没预概形态射，$S$ 连通. 则 $S'$ 连通，且选 $S'$ 中的几何点 $s'$、记 $s$ 为它在 $S$ 中的像，同态
\[
\pi_1(S', s') \longrightarrow \pi_1(S, s)
\]
满射. 若 $g$ 是对预概形的平展覆盖纤维范畴的有效下降态射 (参见 \ref{thm:IX.4.12})，引入 $S'' = S' \times_S S'$ 的几何点 $s'' = \mathrm{diag}(s')$，以及两个投影诱导的两个同态
\[
p_1^{*}, p_2^{*} \colon \pi_1(S'', s'') \longrightarrow \pi_1(S', s'),
\]
则 $\pi_1(S, s)$ 同构于在 Galois 型群范畴中这对态射的余核，i.e. 同构于 $\pi_1(S', s')$ 除以由形如 $p_1^{*}(g'') p_2^{*}(g'')^{-1}$ ($g'' \in \pi_1(S'', s'')$) 的元素生成的闭不变子群.
\end{corollary}

由 \ref{cor:IX.3.4} 确实知道函子 $X \mapsto X \times_S S'$ (从 $S$ 上平展覆盖到 $S'$ 上平展覆盖) 忠实全，这等价于基本群同态是满态射 (\cref{prop:V.6.9}). 最后一个断言是描述 \ref{thm:IX.5.1} 的直接推论.

\section{基本正合序列。具相对连通纤维态射的下降}

\begin{theorem}\label{thm:IX.6.1}
设 $S$ 是剩余域 $k$ 的 Artin 环 $A$ 的谱，$\overline{k}$ 是 $k$ 的代数闭包，$X$ 是 $S$-预概形，$X_0 = X \otimes_A \overline{k}$，$\overline{X}_0 = X \otimes_A \overline{k}$，$a$ 是 $X$ 的几何点，$\overline{a}$ 是它在 $X$ 中的像，$b$ 是它在 $S$ 中的像. 假设 $\overline{X}_0$ 在 $k$ 上拟紧且\emph{几何连通} (N.B. 若 $X$ 在 $S$ 上紧合，这意味着 $H^0(X_0, \mathcal{O}_{X_0})$ 是剩余域在 $k$ 上根式的 Artin 局部环). 则典范同态序列
\[
e \longrightarrow \pi_1(\overline{X}_0, \overline{a}) \longrightarrow \pi_1(X, a) \longrightarrow \pi_1(S, b) \longrightarrow e
\]
正合，且有
\[
\pi_1(S, b) \xleftarrow{\sim} \pi_1(\overline{k}, \overline{k}) = \overline{k} \text{ 在 } k \text{ 上的 伽罗瓦群}.
\]
\end{theorem}

由于基本群不随幂零元素去势而改变，可假设 $A = k$，这已使最后一个同构显然. 设 $k'$ 是 $k$ 在 $\overline{k}$ 中的可分闭包，考虑 $X' = X \otimes_k k'$ 和 $a$ 在 $X'$ 中的像 $a'$. 有典范同态序列
\[
e \longrightarrow \pi_1(\overline{X}_0, \overline{a}) \longrightarrow \pi_1(X', a') \longrightarrow \pi_1(S', b') \longrightarrow e
\]
(其中 $S' = \mathrm{Spec}(k')$). 最后，借助图式
\[
\begin{tikzcd}[row sep=large, column sep=large]
S & X \arrow[l] \\
S' \arrow[u] & X' \arrow[l] \arrow[u]
\end{tikzcd}
\]
有这个序列到 $X/k$ 的相应序列的典范同态. 另一方面看到，这个群序列的同态是同构，如 \ref{cor:IX.4.11} 所示. 故归结到证明第二个序列正合，i.e. 可假设 $k$ 完备. 设 $k_i$ 是 $k$ 在 $\overline{k}$ 中的有限 Galois 子扩张，令 $X_i = X \otimes_k k_i$，设 $a_i$ 是 $a$ 在 $X_i$ 中的像. 把验证自然同态
\[
\pi_1(\overline{X}_0, \overline{a}) \xrightarrow{\sim} \varprojlim_i \pi_1(X_i, a_i)
\]
是同构留给读者，它简单意味着 $X$ 的平展覆盖来自某个 $X_i$ 的平展覆盖，且后者在换成某个 $X_j$ ($j \geqslant i$) 的意义下本质唯一. 另一方面，设 $\pi_i$ 是 $k_i$ 在 $k$ 上的 伽罗瓦群，i.e. $S_i = \mathrm{Spec}(k_i)$ 的 $S$-自同构群的对偶群. 由于函子 $S' \mapsto X \times_S S'$ (从 $S$ 的平展覆盖到 $X$ 的平展覆盖) 忠实全 (\cref{cor:IX.3.4})，得出 $\pi_i$ 也同构于 $X$ 的连通主覆盖 $X_i$ 的 $X$-自同构群的对偶群. 故由 V 6.13 有正合序列
\[
e \longrightarrow \pi_1(X_i, a_i) \longrightarrow \pi_1(X, a) \longrightarrow \pi_i \longrightarrow e.
\]
对这些正合序列取 $i$ 的投射极限，得到正合序列 (因为在 Galois 型群范畴中)，它正是 \ref{thm:IX.6.1} 中考虑的序列. 这完成证明.

\ref{thm:IX.6.1} 中右正合性的几何翻译如下：

\begin{corollary}\label{cor:IX.6.2}
用前面记号，设 $X'$ 是 $X$ 的平展覆盖，$\overline{X}'_0$ 是 $\overline{X}_0$ 的对应平展覆盖. 下列条件等价：
\begin{enumerate}[i)]
\item 存在平展于 $S$ 上的 $S'$ 和 $X$-同构 $X' \xrightarrow{\sim} X \times_S S'$ ($S'$ 此时由 \ref{cor:IX.3.4} 在唯一同构意义下确定).
\item $\overline{X}'_0$ 在 $\overline{X}_0$ 上完全分解.
\end{enumerate}
若 $X'$ 连通，这些条件也等价于：
\begin{enumerate}[i)]
\item[(ii bis)] $\overline{X}'_0$ 在 $\overline{X}_0$ 上有截面.
\end{enumerate}
\end{corollary}

(N.B. 这最后一个补充本质；(i) 与 (ii) 的等价只意味着 $\pi_1(S, b)$ 是 $\pi_1(X, a)$ 除以由 $\pi_1(\overline{X}_0, \overline{a})$ 的像生成的闭不变子群所得的商群，而不是除以该像本身). 在上述条件下，称 $X'$ 是 $X$ 的\emph{几何平凡覆盖}\index[term]{jihepingfanfugai@几何平凡覆盖}.

\begin{remark}\label{rem:IX.6.3}
在断言 \ref{thm:IX.6.1} 中不能把 $k$ 换成 $k$ 的任意代数闭扩张，即使已假设 $k$ 代数闭. 换言之，一般不对：若 $X$ 是代数闭域 $k$ 上的连通代数概形，把 $k$ 换成代数闭扩张时它的基本群不变；例如在特征 $p > 0$ 中对 $k$ 上的仿射直线已经错了，因为无穷远点处有"高分歧"现象，蕴含基本群有"连续"结构. 我们将在下一讲看到，若 $X$ 在 $k$ 上紧合，这类现象不会发生. 还将在 $k$ 特征零时用超越方法证明同样结果.
\end{remark}

\begin{corollary}\label{cor:IX.6.4}
假设 $a$ 定位于 $X$ 中在 $k$ 上有理 (或更一般地剩余域在 $k$ 上根式) 的 $x$ 处. 则正合序列 \ref{thm:IX.6.1} 分裂.
\end{corollary}

可假设 $S = \mathrm{Spec}(k)$. 若 $x$ 在 $k$ 上有理，它对应 $X$ 在 $S$ 上的截面 $S \to X$，把 $b$ 变为 $a$，并定义同态 $\pi_1(S, b) \to \pi_1(X, a)$，即所求分裂. 若 $k(x)$ 在 $k$ 上根式，作基扩张 $\mathrm{Spec}(k(x)) \to \mathrm{Spec}(k)$ 归结到前一情形.

\begin{theorem}\label{thm:IX.6.5}
设 $f \colon X \to S$ 是纤维几何连通的紧合满射有限表现态射，$X'$ 是在 $X$ 上有限表现且紧合的预概形，$s$ 是 $S$ 的点，$F = X_s$ 是 $X$ 在 $s$ 处的纤维，$F'_1$ 是 $X'$ 在 $s$ 处的纤维 $F' = X'_s$ 的连通分量. 为使存在 $F'_1$ 在 $X'$ 中的开邻域 $X'_1$，平展 $S$-概形 $S'_1$ 和 $X$-同构 $X'_1 \xrightarrow{\sim} S'_1 \times_S X$，必要且充分条件是 $X'$ 在 $F'_1$ 的诸点处平展于 $X$ 上，且 $F'_1$ 是 $F$ 的几何平凡覆盖.
\end{theorem}

条件的必要性平凡，剩下证明充分性. 容易归结到 $S$ 诺特的情形. 考虑 $f$ 的 Stein 分解 $X \to T \to S$，其中 $T$ 是 $S$ 上代数 $f_*(\mathcal{O}_X)$ 的谱. 由于 $X$ 在 $S$ 上的诸纤维几何连通，且 $f$ 满射，态射 $T \to S$ 有限满射且根式，故 (\cref{thm:IX.4.10}) 每个平展于 $T$ 上的 $T'$ 来自某个平展于 $S$ 上的 $S'$ 的逆像. 这归结到把 $S$ 换成 $T$ 证明 \ref{thm:IX.6.5}，i.e. 在假设 $f_*(\mathcal{O}_X) = \mathcal{O}_S$ 的情形. 考虑紧合态射 $h \colon X' \to S$ 的 Stein 分解 $X' \to S' \to S$，其中 $S'$ 是代数 $h_*(\mathcal{O}_{X'})$ 的谱. 态射 $X' \to X$ 和 $X' \to S'$ 定义典范态射
\[
X' \longrightarrow X \times_S S',
\]
我们的断言包含在下列断言中：

\begin{corollary}\label{cor:IX.6.6}
设 $f \colon X \to S$ 是局部诺特预概形的紧合态射，使 $f_*(\mathcal{O}_X) = \mathcal{O}_S$，设 $X'$ 是在 $X$ 上紧合的预概形. 考虑 $X' \to S$ 的 Stein 分解 $X' \to S' \to S$ 和典范态射 $X' \to X \times_S S'$. 设 $s$ 是 $S$ 的点，$s'$ 是 $S'$ 的位于 $s$ 上方的点，对应 $X'$ 在 $s$ 处的纤维 $X'_s$ 的连通分量 $F'_1$. 为使态射 $X' \to X \times_S S'$ 在 $s'$ 的平展于 $S$ 上的开邻域 $U'$ 上方是同构，必要且充分条件是 $X'$ 在 $F'_1$ 的诸点处平展于 $X$ 上，且 $F'_1$ 是纤维 $F = X_s$ 的几何平凡覆盖.
\end{corollary}

必要性仍然平凡，剩下证明充分性. 结论也意味着 a) 由 $X' \to X \times_S S'$ 经基变换 $\mathrm{Spec}(\mathcal{O}_{b_{s'}}) \to S'$ 得到的态射是同构，b) $S'$ 在 $s'$ 处平展于 $S$ 上，i.e. $\mathcal{O}_{b_{s'}}$ 在 $\mathcal{O}_{b_s}$ 上平展. 在这种形式下，看到结论在基变换 $\mathrm{Spec}(\mathcal{O}_{b_s}) \to S$ 下不变. 假设也在此基变换下稳定，故可假设 $S$ 是完备诺特局部环的谱. 显然还可假设 $X'$ 连通，此时 $S' = \mathrm{Spec}(\mathcal{O}_{s'})$，$F' = F'_1$.

由于 $X'$ 在 $X$ 上平展的点的集开且包含纤维 $X'_{s'} = F'$，$X'$ 在 $S$ 上紧合，得出 $X'$ 在 $X$ 上平展. 由于它在 $F = X_s$ 上诱导的平展覆盖同构于 $F \otimes_{k(s)} L$，其中 $L$ 在 $k(s)$ 上平展，由 \ref{thm:IX.1.10} 得出它同构于形如 $X \times_S T$ 的覆盖，其中 $T$ 平展于 $S$ 上. (N.B. 这里仍只需用 \ref{thm:IX.1.10} 中函子忠实全的事实，它由 $X$ 上凝聚层的形式同构来自这些层的同构得出). 故若 $T$ 由 $A$ 上有限代数 $B$ 定义，$X'$ 等同于 $X$ 上代数 $\mathcal{O}_X \otimes_A B$ 的谱，由此立即得出 (因为 $f_*(\mathcal{O}_X) = \mathcal{O}_S$) $h_*(\mathcal{O}_{X'})$ 由 $B$ 定义，故典范态射 $X' \to X \times_S S'$ 正是所考虑的同构 $X' \xrightarrow{\sim} X \times_S T$. 这完成证明.

\begin{corollary}\label{cor:IX.6.7}
在 \ref{thm:IX.6.5} 的条件下，为使存在平展于 $S$ 上的预概形 $S'$ 和 $X$-同构 $X' \xrightarrow{\sim} X \times_S S'$，必要且充分条件是 $X'$ 在 $X$ 上平展，且对每个 $s \in S$，$X'_{s'}$ 是 $X_s$ 的几何平凡覆盖.
\end{corollary}

确实，若如此，$X'$ 是同构于平展于 $S$ 上的 $S'_i$ 的逆像的开集 $X'_i$ 的并. 容易看到这些 $S'_i$ 粘合成平展于 $S$ 上的 $S'$，并得到同构 $X' \xrightarrow{\sim} X \times_S S'$. 例如可以说诸 $X'_i$ 配备相对 $X \to S$ 的下降数据，它们必然粘合成整个 $X'$ 上相对 $X \to S$ 的下降数据. 由于后者在诸 $X'_i$ 上有效，容易得出 (借助第 4 号中遗漏的一个推论) 它有效. 也可以把 \ref{cor:IX.6.7} 表述成下列形式：

\begin{corollary}\label{cor:IX.6.8}
设 $f \colon X \to S$ 是纤维几何连通的紧合满射有限表现态射. 则 $f$ 对互相平展有限的预概形的纤维范畴是有效下降态射. 函子 $S' \mapsto X \times_S S'$ 诱导 $S$ 上平展有限预概形范畴与 $X$ 上对每个纤维 $X_s$ 诱导几何平凡覆盖的平展有限预概形范畴的等价.
\end{corollary}

\begin{remark}\label{rem:IX.6.9}
设 $f \colon X \to S$ 是紧合满射态射，$Y$ 局部诺特，则 $f$ 分解为满足 \ref{cor:IX.6.8} 假设的态射 $X \to S'$ 和服从 \ref{thm:IX.4.7} 的有限满射态射 $S' \to S$ 的复合，故为两个对互相平展有限的预概形的纤维范畴都是万有有效下降态射的态射的积. 可以由此得出 $f$ 本身对所述纤维范畴是万有有效下降态射. 于是用不同方法重新得到 4.12.
\end{remark}

\begin{remark}\label{rem:IX.6.10}
把 $f$ 紧合的假设换成：$X$ 在 $S$ 上有限型且在 $S$ 上有截面 (故 $f$ 万有浸没，且是对互相平展预概形的纤维范畴的下降态射)，\ref{cor:IX.6.7} 的结论不再成立，即使 $S$ 是离散赋值环的谱且 $X'$ 是 $X$ 的平展覆盖. 为看到这一点，从 $S$ 上紧合、一般纤维为非奇异有理曲线、特殊纤维 $Z_0$ 由两条共点直线组成的 $Z$ 出发. 例如，若 $t$ 是赋值环 $A$ 的幺模化元，取 $\mathbf{P}^2_A$ 的由齐次方程 $x^2 + y^2 + tz^2 = 0$ (齐次坐标 $x, y, z$) 定义的闭子概形 $Z$. 取 $X$ 为 $Z_0$ 的奇异点 $a$ 在并 $Z \cup \mathbf{P}^2_k$ 中的余集. $X$ 的诸纤维是 $\mathbf{P}^1_k$ 和 $\mathbf{P}^2_k - a$，故几何单连通 (i.e. 这样的纤维的每个平展覆盖几何平凡). 然而，如第 4 号那样进行，容易构造 $X$ 的平展覆盖，它不来自 $S$ 的平展覆盖，即把 $Z - a$ 和 $\mathbf{P}^2_k - a$ 的平凡覆盖粘合得到. 反之，把紧合性假设换成 $X$ 在 $S$ 上万有开有限表现，\ref{cor:IX.6.7} 的结论可能仍然成立\footnote{现在已证明，$g$ 只需万有开且满射；参见 SGA 4 XV 1.15.}. 至少假设 $X$ 在 $S$ 上的诸纤维几何不可约 (而不只是几何连通) 时如此. 只指出在这个问题中可以归结到 $S$ 是剩余域代数闭的完备离散赋值环的谱的情形.
\end{remark}

\ref{cor:IX.6.7} 在基本群术语中的解释如下：

\begin{corollary}\label{cor:IX.6.11}
设 $f \colon X \to S$ 是纤维几何连通的紧合满射有限表现态射. 假设 $X$ 故 $S$ 连通. 设 $a$ 是 $X$ 的几何点，$b$ 是它在 $S$ 中的像，对每个 $s \in S$，选 $k(s)$ 的代数闭包 $\overline{k(s)}$，在此扩张中取值的 $S_s$ 的几何点 $a_s$，以及从 $a_s$ 到 $a$ 的道路类，从而得到同态 $\pi_1(\overline{X}_s, a_s) \to \pi_1(X, a)$，其中 $\overline{X}_s = X_s \otimes_{k(s)} \overline{k(s)}$. 则同态 $\pi_1(X, a) \to \pi_1(S, b)$ 满射，且它的核是 $\pi_1(X, a)$ 的由诸 $\pi_1(\overline{X}_s, a_s)$ 的像生成的闭不变子群.
\end{corollary}

\begin{remark}\label{rem:IX.6.12}
在 \ref{cor:IX.6.7} 的条件下，假设 $S$ 诺特，容易看到使 $S'_{s'}$ 在 $X_s$ 上几何平凡的点 $s \in S$ 的集是可构集；若 $X'$ 在 $X$ 上紧合，它甚至开，如 \ref{cor:IX.6.6} 所见. 故若 $S$ 是 Jacobson 预概形 (例如域上有限型)，或 $X'$ 在 $X$ 上紧合，验证 \ref{cor:IX.6.7} 的条件时可限于 $S$ 的闭点 $s$. 同样，在 \ref{cor:IX.6.11} 中此时只需取 $S$ 的闭点的诸 $\pi_1(\overline{X}_s, a_s)$.
\end{remark}

\section{文献}\printbibliography[heading=none]


